% GENERATED FILE. Edit canonical lessons, then run npm run generate:books.
% canonical-source-hash: fnv1a64:765c9f1d161ca73b
% canonical-lessons: GE-C453-klippe, GE-C453-hochebene, GE-C453-gipfel, GE-C453-abhang, GE-C453-schlucht, GE-R453-review-446-453

\chapter{Berge und Täler}
\label{ch:ge-berge-und-taler}
\hlchaptermodality{453}

\begin{chapteropening}
\textbf{By the end of this chapter:} \emph{I can describe high and steep ground: the cliff, the plateau, the summit, the slope and the gorge.}

Complete the last lesson of chapter 453: I can describe high and steep ground: the cliff, the plateau, the summit, the slope and the gorge.
\end{chapteropening}

\section[die Klippe]{die Klippe — the cliff; the rock (in the sea)}

\label{lesson:GE-C453-klippe}



\begin{quote}
Before the new one: say the German for to deliver, then the German for to exchange.
\end{quote}

\subsection*{You'll want to know: die Klippe}
\textbf{die Klippe} — \textquotedblleft{}the cliff; the rock (in the sea)\textquotedblright{}.

\begin{quote}
Von der Klippe sieht man weit über das Meer.
\end{quote}

\begin{quote}
\emph{From the cliff you can see far out over the sea.}
\end{quote}

Plural: \emph{die Klippen}.

\begin{cousinweb}
\emph{Klippe} goes back to medieval German \emph{klippe} 'rock, cliff' (EtymDB) and came into the standard language from Low German; Dutch \emph{klip} 'cliff' is a cousin.
\end{cousinweb}

\begin{grammarlens}[title={What you've built}]
One more word for telling, explaining and describing at B1.
\end{grammarlens}

\subsection*{Your turn}
\begin{itemize}
\raggedright
  \item \emph{Say it:} \emph{die Klippe}
  \item \emph{Say it:} \emph{die Klippe}, once more
  \item \emph{Say it:} a sentence of your own with \emph{die Klippe}
  \item \emph{Recall:} say the German for to deliver, then the German for to exchange, then say \emph{die Klippe} again
\end{itemize}

\begin{tcolorbox}[breakable,colback=teal!4,colframe=teal!35!black,title={Before you move on}]
What does die Klippe mean? (\textquotedblleft{}The cliff; the rock (in the sea)\textquotedblright{}.) Say it once more.
\end{tcolorbox}
\section[die Hochebene]{die Hochebene — the plateau}

\label{lesson:GE-C453-hochebene}



\begin{quote}
Before the new one: say the German for to exchange, then the German for the cliff.
\end{quote}

\subsection*{You'll want to know: die Hochebene}
\textbf{die Hochebene} — \textquotedblleft{}the plateau\textquotedblright{}.

\begin{quote}
Auf der Hochebene ist es kühl und windig.
\end{quote}

\begin{quote}
\emph{On the plateau it is cool and windy.}
\end{quote}

Plural: \emph{die Hochebenen}.

\begin{cousinweb}
\emph{Hochebene} is built from \emph{hoch} 'high', inherited from Proto-Germanic \emph{hauhaz} (EtymDB), and \emph{Ebene} 'plain', from Old High German \emph{ebanī} 'flatness' (EtymDB).
\end{cousinweb}

\begin{grammarlens}[title={What you've built}]
One more word for telling, explaining and describing at B1.
\end{grammarlens}

\subsection*{Your turn}
\begin{itemize}
\raggedright
  \item \emph{Say it:} \emph{die Hochebene}
  \item \emph{Say it:} \emph{die Hochebene}, once more
  \item \emph{Say it:} a sentence of your own with \emph{die Hochebene}
  \item \emph{Recall:} say the German for to exchange, then the German for the cliff, then say \emph{die Hochebene} again
\end{itemize}

\begin{tcolorbox}[breakable,colback=teal!4,colframe=teal!35!black,title={Before you move on}]
What does die Hochebene mean? (\textquotedblleft{}The plateau\textquotedblright{}.) Say it once more.
\end{tcolorbox}
\section[der Gipfel]{der Gipfel — the summit, the peak; a summit meeting}

\label{lesson:GE-C453-gipfel}



\begin{quote}
Before the new one: say the German for the cliff, then the German for the plateau.
\end{quote}

\subsection*{You'll want to know: der Gipfel}
\textbf{der Gipfel} — \textquotedblleft{}the summit, the peak; a summit meeting\textquotedblright{}.

\begin{quote}
Nach vier Stunden sind wir auf dem Gipfel.
\end{quote}

\begin{quote}
\emph{After four hours we are at the summit.}
\end{quote}

Same form in the plural. Also a political summit.

\begin{cousinweb}
The origin of \emph{Gipfel} is not settled. EtymDB lists the word with the sense 'summit' but gives no source.
\end{cousinweb}

\begin{grammarlens}[title={What you've built}]
One more word for telling, explaining and describing at B1.
\end{grammarlens}

\subsection*{Your turn}
\begin{itemize}
\raggedright
  \item \emph{Say it:} \emph{der Gipfel}
  \item \emph{Say it:} \emph{der Gipfel}, once more
  \item \emph{Say it:} a sentence of your own with \emph{der Gipfel}
  \item \emph{Recall:} say the German for the cliff, then the German for the plateau, then say \emph{der Gipfel} again
\end{itemize}

\begin{tcolorbox}[breakable,colback=teal!4,colframe=teal!35!black,title={Before you move on}]
What does der Gipfel mean? (\textquotedblleft{}The summit, the peak; a summit meeting\textquotedblright{}.) Say it once more.
\end{tcolorbox}
\section[der Abhang]{der Abhang — the slope, the hillside}

\label{lesson:GE-C453-abhang}



\begin{quote}
Before the new one: say the German for the plateau, then the German for the summit.
\end{quote}

\subsection*{You'll want to know: der Abhang}
\textbf{der Abhang} — \textquotedblleft{}the slope, the hillside\textquotedblright{}.

\begin{quote}
Am steilen Abhang wachsen nur kleine Bäume.
\end{quote}

\begin{quote}
\emph{Only small trees grow on the steep slope.}
\end{quote}

Plural: \emph{die Abhänge}. Also \emph{der Hang}.

\begin{cousinweb}
\emph{Abhang} is built from \emph{ab} 'down' and \emph{Hang} 'slope', related to \emph{hängen} 'to hang' (EtymDB lists \emph{Abhang}): ground that hangs downwards.
\end{cousinweb}

\begin{grammarlens}[title={What you've built}]
One more word for telling, explaining and describing at B1.
\end{grammarlens}

\subsection*{Your turn}
\begin{itemize}
\raggedright
  \item \emph{Say it:} \emph{der Abhang}
  \item \emph{Say it:} \emph{der Abhang}, once more
  \item \emph{Say it:} a sentence of your own with \emph{der Abhang}
  \item \emph{Recall:} say the German for the plateau, then the German for the summit, then say \emph{der Abhang} again
\end{itemize}

\begin{tcolorbox}[breakable,colback=teal!4,colframe=teal!35!black,title={Before you move on}]
What does der Abhang mean? (\textquotedblleft{}The slope, the hillside\textquotedblright{}.) Say it once more.
\end{tcolorbox}
\section[die Schlucht]{die Schlucht — the gorge, the ravine}

\label{lesson:GE-C453-schlucht}



\begin{quote}
Before the new one: say the German for the summit, then the German for the slope.
\end{quote}

\subsection*{You'll want to know: die Schlucht}
\textbf{die Schlucht} — \textquotedblleft{}the gorge, the ravine\textquotedblright{}.

\begin{quote}
Unten in der tiefen Schlucht ist ein kleiner Fluss.
\end{quote}

\begin{quote}
\emph{Down in the deep gorge there is a small river.}
\end{quote}

Plural: \emph{die Schluchten}.

\begin{cousinweb}
\emph{Schlucht} came into the standard language from Low German; it is usually linked to \emph{schlucken} 'to swallow', and EtymDB lists Old Norse \emph{slok} 'gully, channel' as a cousin.
\end{cousinweb}

\begin{grammarlens}[title={What you've built}]
One more word for telling, explaining and describing at B1.
\end{grammarlens}

\subsection*{Your turn}
\begin{itemize}
\raggedright
  \item \emph{Say it:} \emph{die Schlucht}
  \item \emph{Say it:} \emph{die Schlucht}, once more
  \item \emph{Say it:} a sentence of your own with \emph{die Schlucht}
  \item \emph{Recall:} say the German for the summit, then the German for the slope, then say \emph{die Schlucht} again
\end{itemize}

\begin{tcolorbox}[breakable,colback=teal!4,colframe=teal!35!black,title={Before you move on}]
What does die Schlucht mean? (\textquotedblleft{}The gorge, the ravine\textquotedblright{}.) Say it once more.
\end{tcolorbox}
\section[(dialogue)]{Words from chapters 446–453 — a review}

\label{lesson:GE-R453-review-446-453}



\begin{quote}
Say the German for the dry cleaner and the launderette.
\end{quote}

\subsection*{Your turn}
Say these aloud:

\begin{itemize}
\raggedright
  \item the dry cleaner, the launderette, to remove, to patch, to dry
  \item the hairdresser, to blow-dry, the optician, the watchmaker, the jeweller
  \item to put postage on, to send, to post, the small parcel, the registered letter
  \item the cashier, to take the money, the cash, the card reader, the change
  \item the cash machine, the PIN, the bank statement, to redeem, to pay out
  \item to reduce, the special offer, the bargain, to haggle, the flea market
  \item the packaging, to wrap up, the cardboard box, to deliver, to exchange
  \item the cliff, the plateau, the summit, the slope, the gorge
\end{itemize}

\begin{tcolorbox}[breakable,colback=teal!4,colframe=teal!35!black,title={Before you move on}]
The dry cleaner? (\textbf{die Reinigung}.) The launderette? (\textbf{der Waschsalon}.) To remove? (\textbf{entfernen}.) To patch? (\textbf{flicken}.) To dry? (\textbf{trocknen}.) The hairdresser? (\textbf{der Friseur}.) To blow-dry? (\textbf{föhnen}.) The optician? (\textbf{der Optiker}.) The watchmaker? (\textbf{der Uhrmacher}.) The jeweller? (\textbf{der Juwelier}.) To put postage on? (\textbf{frankieren}.) To send? (\textbf{verschicken}.) To post? (\textbf{einwerfen}.) The small parcel? (\textbf{das Päckchen}.) The registered letter? (\textbf{das Einschreiben}.) The cashier? (\textbf{der Kassierer}.) To take the money? (\textbf{kassieren}.) The cash? (\textbf{das Bargeld}.) The card reader? (\textbf{das Kartenlesegerät}.) The change? (\textbf{das Wechselgeld}.) The cash machine? (\textbf{der Geldautomat}.) The PIN? (\textbf{die Geheimzahl}.) The bank statement? (\textbf{der Kontoauszug}.) To redeem? (\textbf{einlösen}.) To pay out? (\textbf{auszahlen}.) To reduce? (\textbf{reduzieren}.) The special offer? (\textbf{das Sonderangebot}.) The bargain? (\textbf{das Schnäppchen}.) To haggle? (\textbf{feilschen}.) The flea market? (\textbf{der Flohmarkt}.) The packaging? (\textbf{die Verpackung}.) To wrap up? (\textbf{verpacken}.) The cardboard box? (\textbf{der Karton}.) To deliver? (\textbf{liefern}.) To exchange? (\textbf{umtauschen}.) The cliff? (\textbf{die Klippe}.) The plateau? (\textbf{die Hochebene}.) The summit? (\textbf{der Gipfel}.) The slope? (\textbf{der Abhang}.) The gorge? (\textbf{die Schlucht}.)
\end{tcolorbox}
